\section{Non-Functional Evidence, Testing and DevOps}
\label{sec:nfr-testing-devops}

\subsection{Performance and Load Testing}
The defined workload is 50 concurrent virtual users against a read-heavy mix
of dashboard, list and detail endpoints for five minutes after a warm-up
period, using k6. The proposed acceptance targets are a 95th-percentile
response time under 500\,ms for read endpoints and under 1000\,ms for write
endpoints, with an unexpected-error rate below 1\%. Login and report-export
endpoints are measured separately from the main workload, since they are
deliberately rate-limited at the gateway. \cref{tab:load} is the template used
to record the actual run; it should be completed with real numbers and
hardware details from \texttt{tests/load/} before final submission.

\begin{table}[H]
  \centering
  \caption{Load-test result template (complete with measured values)}
  \label{tab:load}
  \begin{tabular}{L{3.4cm}R{2.2cm}R{2.2cm}R{2.2cm}R{2.2cm}}
    \toprule
    \rowcolor{ceTeal}\thd{Workload} & \thd{VUs} & \thd{p95 (ms)} & \thd{Errors (\%)} & \thd{Duration} \\
    \midrule
    Dashboard reads & 50 & -- & -- & 5 min \\
    Registration write & 10 & -- & -- & 2 min \\
    Payment intent + callback & 10 & -- & -- & 2 min \\
    \bottomrule
  \end{tabular}
\end{table}

\subsection{Scalability Strategy}
API instances are stateless with respect to authentication (JWT verification
uses only the public key, held by every instance), so the gateway can load
balance multiple replicas of any domain service without sticky sessions.
Identity's session and refresh-token state lives centrally in
\texttt{identity\_db}, which remains the single source of truth for session
revocation across replicas. Consumers use a bounded prefetch count so that
adding worker replicas increases throughput without unbounded memory growth,
and every consumer transaction is idempotent, so redelivery after a crash or
a rebalance cannot duplicate side effects. \cref{fig:d5-deployment}'s
Compose topology is the current single-node profile; horizontally scaling any
one service is, by this design, a configuration change (replica count) rather
than a code change.

\subsection{Testing and Quality Assurance}
Each service owns a \texttt{tests/} suite of unit tests (pure domain-function
logic, e.g. conflict detection, at-risk evaluation, payroll calculation) and
integration tests (through the FastAPI test client, covering full HTTP
request/response cycles, tenant isolation, and concurrency/rollback
scenarios). Coverage is measured on each service's \texttt{app.domain}
package specifically, so that trivial routing code cannot inflate the
percentage; the CI pipeline fails the build below 70\%. \cref{tab:coverage}
reports the coverage achieved at the time of writing (see
\texttt{docs/evidence/coverage.txt} for the full \texttt{pytest-cov} report).

\begin{table}[H]
  \centering\small
  \caption{Automated test coverage by service (core domain logic)}
  \label{tab:coverage}
  \begin{tabular}{L{2.6cm}R{2.0cm}R{2.0cm}R{2.0cm}}
    \toprule
    \rowcolor{ceTeal}\thd{Service} & \thd{Tests} & \thd{Domain coverage} & \thd{Status} \\
    \midrule
    Identity & 15 & 96\% & Passing \\
    Academic & 21 & $\sim$89\% & Passing \\
    Finance & 8 & $\sim$88\% & Passing \\
    HR & 11 & $\sim$88\% & Passing \\
    \bottomrule
  \end{tabular}
\end{table}

A written test plan maps each functional requirement (Appendix~\ref{app:trace})
to the automated test that verifies it. Open and resolved issues are tracked
on a GitHub Issues board (linked from the repository README); screenshots or
an export of that board should be attached as evidence before submission.

\subsection{DevOps: Containerisation and CI/CD}
Every service (and its worker) has its own multi-stage Dockerfile, running as
a non-root user, with a container \texttt{HEALTHCHECK} against its
\texttt{/health/ready} endpoint. \texttt{docker-compose.yml} brings up the
full system -- gateway, four services, three workers, the mobile-money
simulator, PostgreSQL, RabbitMQ, Mailpit and the observability stack -- with
one command, gated by explicit \texttt{depends\_on} health and
completion conditions so that migrations run before the API starts and the
API is healthy before the gateway routes to it. The GitHub Actions pipeline
(\texttt{.github/workflows/ci.yml}) runs on every pull request and on every
push to \texttt{main}: it lints and tests each Python service independently
with a coverage gate, builds and type-checks the frontend, and -- on
\texttt{main} only -- builds every service's container image, verifying the
Dockerfiles are still valid without publishing to a registry.

\subsection{Deployment Guide (Summary)}
A third party can run the system from a clean machine with Docker Desktop
installed by cloning the repository, running
\texttt{python scripts/bootstrap.py} (generates \texttt{.env} and local JWT
signing keys), then \texttt{docker compose up --build}. The full step-by-step
guide, including the non-Docker quick-start (\texttt{run.bat}) used for this
report's screenshots, is kept in the repository README and is not duplicated
here to preserve the page budget.
